% Options for packages loaded elsewhere
\PassOptionsToPackage{unicode}{hyperref}
\PassOptionsToPackage{hyphens}{url}
%
\documentclass[
]{article}
\usepackage{amsmath,amssymb}
\usepackage{iftex}
\ifPDFTeX
  \usepackage[T1]{fontenc}
  \usepackage[utf8]{inputenc}
  \usepackage{textcomp} % provide euro and other symbols
\else % if luatex or xetex
  \usepackage{unicode-math} % this also loads fontspec
  \defaultfontfeatures{Scale=MatchLowercase}
  \defaultfontfeatures[\rmfamily]{Ligatures=TeX,Scale=1}
\fi
\usepackage{lmodern}
\ifPDFTeX\else
  % xetex/luatex font selection
\fi
% Use upquote if available, for straight quotes in verbatim environments
\IfFileExists{upquote.sty}{\usepackage{upquote}}{}
\IfFileExists{microtype.sty}{% use microtype if available
  \usepackage[]{microtype}
  \UseMicrotypeSet[protrusion]{basicmath} % disable protrusion for tt fonts
}{}
\makeatletter
\@ifundefined{KOMAClassName}{% if non-KOMA class
  \IfFileExists{parskip.sty}{%
    \usepackage{parskip}
  }{% else
    \setlength{\parindent}{0pt}
    \setlength{\parskip}{6pt plus 2pt minus 1pt}}
}{% if KOMA class
  \KOMAoptions{parskip=half}}
\makeatother
\usepackage{xcolor}
\usepackage{color}
\usepackage{fancyvrb}
\newcommand{\VerbBar}{|}
\newcommand{\VERB}{\Verb[commandchars=\\\{\}]}
\DefineVerbatimEnvironment{Highlighting}{Verbatim}{commandchars=\\\{\}}
% Add ',fontsize=\small' for more characters per line
\newenvironment{Shaded}{}{}
\newcommand{\AlertTok}[1]{\textcolor[rgb]{1.00,0.00,0.00}{\textbf{#1}}}
\newcommand{\AnnotationTok}[1]{\textcolor[rgb]{0.38,0.63,0.69}{\textbf{\textit{#1}}}}
\newcommand{\AttributeTok}[1]{\textcolor[rgb]{0.49,0.56,0.16}{#1}}
\newcommand{\BaseNTok}[1]{\textcolor[rgb]{0.25,0.63,0.44}{#1}}
\newcommand{\BuiltInTok}[1]{\textcolor[rgb]{0.00,0.50,0.00}{#1}}
\newcommand{\CharTok}[1]{\textcolor[rgb]{0.25,0.44,0.63}{#1}}
\newcommand{\CommentTok}[1]{\textcolor[rgb]{0.38,0.63,0.69}{\textit{#1}}}
\newcommand{\CommentVarTok}[1]{\textcolor[rgb]{0.38,0.63,0.69}{\textbf{\textit{#1}}}}
\newcommand{\ConstantTok}[1]{\textcolor[rgb]{0.53,0.00,0.00}{#1}}
\newcommand{\ControlFlowTok}[1]{\textcolor[rgb]{0.00,0.44,0.13}{\textbf{#1}}}
\newcommand{\DataTypeTok}[1]{\textcolor[rgb]{0.56,0.13,0.00}{#1}}
\newcommand{\DecValTok}[1]{\textcolor[rgb]{0.25,0.63,0.44}{#1}}
\newcommand{\DocumentationTok}[1]{\textcolor[rgb]{0.73,0.13,0.13}{\textit{#1}}}
\newcommand{\ErrorTok}[1]{\textcolor[rgb]{1.00,0.00,0.00}{\textbf{#1}}}
\newcommand{\ExtensionTok}[1]{#1}
\newcommand{\FloatTok}[1]{\textcolor[rgb]{0.25,0.63,0.44}{#1}}
\newcommand{\FunctionTok}[1]{\textcolor[rgb]{0.02,0.16,0.49}{#1}}
\newcommand{\ImportTok}[1]{\textcolor[rgb]{0.00,0.50,0.00}{\textbf{#1}}}
\newcommand{\InformationTok}[1]{\textcolor[rgb]{0.38,0.63,0.69}{\textbf{\textit{#1}}}}
\newcommand{\KeywordTok}[1]{\textcolor[rgb]{0.00,0.44,0.13}{\textbf{#1}}}
\newcommand{\NormalTok}[1]{#1}
\newcommand{\OperatorTok}[1]{\textcolor[rgb]{0.40,0.40,0.40}{#1}}
\newcommand{\OtherTok}[1]{\textcolor[rgb]{0.00,0.44,0.13}{#1}}
\newcommand{\PreprocessorTok}[1]{\textcolor[rgb]{0.74,0.48,0.00}{#1}}
\newcommand{\RegionMarkerTok}[1]{#1}
\newcommand{\SpecialCharTok}[1]{\textcolor[rgb]{0.25,0.44,0.63}{#1}}
\newcommand{\SpecialStringTok}[1]{\textcolor[rgb]{0.73,0.40,0.53}{#1}}
\newcommand{\StringTok}[1]{\textcolor[rgb]{0.25,0.44,0.63}{#1}}
\newcommand{\VariableTok}[1]{\textcolor[rgb]{0.10,0.09,0.49}{#1}}
\newcommand{\VerbatimStringTok}[1]{\textcolor[rgb]{0.25,0.44,0.63}{#1}}
\newcommand{\WarningTok}[1]{\textcolor[rgb]{0.38,0.63,0.69}{\textbf{\textit{#1}}}}
\usepackage{longtable,booktabs,array}
\usepackage{calc} % for calculating minipage widths
% Correct order of tables after \paragraph or \subparagraph
\usepackage{etoolbox}
\makeatletter
\patchcmd\longtable{\par}{\if@noskipsec\mbox{}\fi\par}{}{}
\makeatother
% Allow footnotes in longtable head/foot
\IfFileExists{footnotehyper.sty}{\usepackage{footnotehyper}}{\usepackage{footnote}}
\makesavenoteenv{longtable}
\setlength{\emergencystretch}{3em} % prevent overfull lines
\providecommand{\tightlist}{%
  \setlength{\itemsep}{0pt}\setlength{\parskip}{0pt}}
\setcounter{secnumdepth}{-\maxdimen} % remove section numbering
\ifLuaTeX
  \usepackage{selnolig}  % disable illegal ligatures
\fi
\IfFileExists{bookmark.sty}{\usepackage{bookmark}}{\usepackage{hyperref}}
\IfFileExists{xurl.sty}{\usepackage{xurl}}{} % add URL line breaks if available
\urlstyle{same}
\hypersetup{
  hidelinks,
  pdfcreator={LaTeX via pandoc}}

\author{}
\date{}

\begin{document}

\hypertarget{satspath-resolver-semantics-v1}{%
\section{SatsPath Resolver Semantics
v1}\label{satspath-resolver-semantics-v1}}

Resolver provenance and verification are reported separately. A
transparent result carries the signed profile, latest name event,
inclusion proof, signed checkpoint, optional consistency proof and
identifier attestation. Implementations expose independent states for
profile signature, identifier binding, key continuity, log inclusion,
checkpoint consistency and payment-method ownership; these must not be
collapsed into a single \texttt{verified} boolean.

SatsPath is transport-neutral. A resolver is any component that can map
a SatsPath identifier to a \texttt{SignedPaymentProfile}.

Resolvers do not decide whether money moves. Resolvers only discover
signed profile data. The quote pipeline verifies and routes after
resolution.

\hypertarget{resolver-interface}{%
\subsection{Resolver Interface}\label{resolver-interface}}

Conceptually, every resolver implements:

\begin{Shaded}
\begin{Highlighting}[]
\NormalTok{resolve(identifier) {-}\textgreater{} SignedPaymentProfile | NotFound | Unavailable | Invalid}
\end{Highlighting}
\end{Shaded}

The Rust reference implementation exposes this as
\texttt{ProfileResolver}:

\begin{Shaded}
\begin{Highlighting}[]
\KeywordTok{async} \KeywordTok{fn}\NormalTok{ resolve\_alias(}\OperatorTok{\&}\KeywordTok{self}\OperatorTok{,}\NormalTok{ alias}\OperatorTok{:} \OperatorTok{\&}\DataTypeTok{str}\NormalTok{) }\OperatorTok{{-}\textgreater{}} \DataTypeTok{Result}\OperatorTok{\textless{}}\NormalTok{SignedPaymentProfile}\OperatorTok{\textgreater{};}
\end{Highlighting}
\end{Shaded}

\texttt{AliasNotFound} maps to \texttt{NotFound}. Network or transport
errors map to \texttt{Unavailable} unless the resolver can prove the
returned data is authoritative and invalid.

\hypertarget{resolver-chain}{%
\subsection{Resolver Chain}\label{resolver-chain}}

The default resolver chain is:

\begin{Shaded}
\begin{Highlighting}[]
\NormalTok{local registry {-}\textgreater{} BIP{-}353 {-}\textgreater{} HTTPS {-}\textgreater{} Nostr}
\end{Highlighting}
\end{Shaded}

Implementations MAY add other resolvers, including P2P transports,
platform APIs, QR-imported profiles, or hardware-wallet-provided profile
stores.

The chain MUST:

\begin{enumerate}
\def\labelenumi{\arabic{enumi}.}
\tightlist
\item
  Normalize the identifier.
\item
  Query each resolver in order.
\item
  Continue past \texttt{NotFound}.
\item
  Continue past transient \texttt{Unavailable} failures.
\item
  Verify any returned signed profile before routing.
\item
  Stop at the first verified profile.
\item
  Return \texttt{not\_registered} if no resolver returns a usable
  profile.
\end{enumerate}

The chain MUST NOT trust a resolver because of its transport. An HTTPS
response, DNS response, Nostr event, Pear/Holepunch message, or local
file all require the same signature verification.

\hypertarget{identifier-verification}{%
\subsection{Identifier Verification}\label{identifier-verification}}

SatsPath separates profile verification from identifier verification:

\begin{itemize}
\tightlist
\item
  Profile verification means the \texttt{SignedPaymentProfile} signature
  verifies against \texttt{profile.identity\_pubkey}.
\item
  Identifier verification means an external authority or challenge
  proves that the canonical identifier was controlled or approved for
  that profile.
\end{itemize}

Plain email syntax, such as \texttt{alice@gmail.com}, is only an
identifier. It does not prove inbox access, domain ownership, or
payment-method ownership.

DNSSEC/BIP-353 and Nostr/NIP-05 can verify publication for domains the
receiver controls. Consumer email domains usually cannot publish records
under their own domain, so those identifiers require an explicit
platform challenge resolver or fall back to the invite flow.

The daemon sets \texttt{identifier\_verified} only when an attestation
binds the exact identifier hash, key and profile hash in the latest
event and its verifier key and method are explicitly trusted. Otherwise
it remains false and first contact is TOFU.

\hypertarget{local-registry-resolver}{%
\subsection{Local Registry Resolver}\label{local-registry-resolver}}

The local registry is a file-backed resolver used by CLI and daemon
flows.

Reference files:

\begin{itemize}
\tightlist
\item
  \texttt{.satspath/satspath-transparency-v1.sqlite3} (daemon security
  state)
\item
  \texttt{.satspath/registry.json} (legacy CLI/local resolver
  compatibility only)
\item
  \texttt{.satspath/peers/registry.local.json}
\end{itemize}

The local registry MAY store profiles created locally or imported from
another transport. A profile imported into the registry MUST be verified
before use.

\hypertarget{bip-353-dns-resolver}{%
\subsection{BIP-353 DNS Resolver}\label{bip-353-dns-resolver}}

BIP-353 resolves a Bitcoin payment instruction from DNS TXT records. In
SatsPath v1, BIP-353 can provide:

\begin{itemize}
\tightlist
\item
  A direct \texttt{bitcoin:} / BIP-321 payment instruction.
\item
  A pointer to a SatsPath signed profile, such as \texttt{sp-profile}.
\item
  A hash commitment, such as \texttt{sp-profile-hash}.
\end{itemize}

Strict DNS behavior:

\begin{itemize}
\tightlist
\item
  DNSSEC validation is required for trustworthy mainnet use.
\item
  If no DNSSEC-validating resolver is available, the implementation MUST
  fail closed.
\item
  Development mode MAY accept unvalidated DNS only behind an explicit
  flag.
\end{itemize}

The reference implementation provides:

\begin{itemize}
\tightlist
\item
  \texttt{DohTxtResolver}
\item
  \texttt{resolve\_bip353\_with}
\item
  \texttt{DnssecPolicy::Strict}
\item
  \texttt{DnssecPolicy::DevInsecure}
\end{itemize}

\hypertarget{https-resolver}{%
\subsection{HTTPS Resolver}\label{https-resolver}}

The HTTPS resolver maps an identifier to a well-known signed profile
endpoint.

Expected endpoint shape:

\begin{Shaded}
\begin{Highlighting}[]
\NormalTok{https://\textless{}domain\textgreater{}/.well{-}known/satspath/\textless{}user\textgreater{}}
\end{Highlighting}
\end{Shaded}

The resolver MUST:

\begin{itemize}
\tightlist
\item
  Fetch JSON over HTTPS.
\item
  Decode as \texttt{SignedPaymentProfile}.
\item
  Verify the profile signature.
\item
  Reject malformed JSON.
\item
  Treat HTTP 404 as \texttt{NotFound}.
\item
  Treat transport failures as \texttt{Unavailable}.
\end{itemize}

HTTP transport security does not replace profile signatures.

\hypertarget{nostr-resolver}{%
\subsection{Nostr Resolver}\label{nostr-resolver}}

The Nostr resolver is a transport for discovering signed profile data
through Nostr events. It does not require Pear/Holepunch or any
SatsPath-specific P2P network.

Resolution flow:

\begin{enumerate}
\def\labelenumi{\arabic{enumi}.}
\tightlist
\item
  Parse \texttt{user@domain}.
\item
  Fetch \texttt{https://domain/.well-known/nostr.json?name=user}.
\item
  Read the NIP-05 pubkey and relay hints.
\item
  Query relays for a Nostr kind \texttt{30078} event authored by that
  pubkey.
\item
  Require a \texttt{d} tag equal to
  \texttt{satspath-profile:\textless{}canonical\ user@domain\textgreater{}}.
\item
  Parse the event \texttt{content} as \texttt{SignedPaymentProfile}.
\item
  Verify the SatsPath profile signature and expiry.
\end{enumerate}

Nostr event signatures and SatsPath profile signatures are separate
layers. The Nostr event proves which Nostr key published the event. The
SatsPath signature proves that the payment profile is controlled by the
SatsPath protocol identity key.

The resolver MUST still return a \texttt{SignedPaymentProfile} and the
quote pipeline MUST verify the SatsPath profile signature.

NIP-05 verifies that the domain's well-known metadata maps the local
identifier part to a Nostr pubkey. It does not verify a consumer email
inbox unless that same domain/account infrastructure explicitly performs
such verification.

Reference event shape:

\begin{Shaded}
\begin{Highlighting}[]
\FunctionTok{\{}
  \DataTypeTok{"kind"}\FunctionTok{:} \DecValTok{30078}\FunctionTok{,}
  \DataTypeTok{"pubkey"}\FunctionTok{:} \StringTok{"\textless{}nip05 nostr pubkey hex\textgreater{}"}\FunctionTok{,}
  \DataTypeTok{"tags"}\FunctionTok{:} \OtherTok{[[}\StringTok{"d"}\OtherTok{,} \StringTok{"satspath{-}profile:alice@example.com"}\OtherTok{]]}\FunctionTok{,}
  \DataTypeTok{"content"}\FunctionTok{:} \StringTok{"\{}\CharTok{\textbackslash{}"}\StringTok{profile}\CharTok{\textbackslash{}"}\StringTok{:\{...\},}\CharTok{\textbackslash{}"}\StringTok{signature}\CharTok{\textbackslash{}"}\StringTok{:}\CharTok{\textbackslash{}"}\StringTok{3044...}\CharTok{\textbackslash{}"}\StringTok{\}"}
\FunctionTok{\}}
\end{Highlighting}
\end{Shaded}

Implementations MAY configure fallback relays. The Rust implementation
reads \texttt{SATSPATH\_NOSTR\_RELAYS} as a comma-separated list and
otherwise uses default public relays.

\hypertarget{p2p-resolver}{%
\subsection{P2P Resolver}\label{p2p-resolver}}

Pear/Holepunch P2P transport is optional. It can announce, request, and
return signed profiles between peers.

P2P discovery MUST NOT be treated as weaker or stronger than HTTPS or
DNS by default. The returned profile is still untrusted until:

\begin{enumerate}
\def\labelenumi{\arabic{enumi}.}
\tightlist
\item
  The response decodes as \texttt{SignedPaymentProfile}.
\item
  The SatsPath profile signature verifies.
\item
  Expiry checks pass.
\item
  Optional method ownership checks pass.
\end{enumerate}

P2P wire behavior is specified in \href{./wire_p2p.md}{wire\_p2p.md}.

\hypertarget{resolver-output-and-quote-mapping}{%
\subsection{Resolver Output and Quote
Mapping}\label{resolver-output-and-quote-mapping}}

Resolver outcome maps to quote behavior as follows:

\begin{longtable}[]{@{}
  >{\raggedright\arraybackslash}p{(\columnwidth - 2\tabcolsep) * \real{0.5083}}
  >{\raggedright\arraybackslash}p{(\columnwidth - 2\tabcolsep) * \real{0.4917}}@{}}
\toprule\noalign{}
\begin{minipage}[b]{\linewidth}\raggedright
Resolver outcome
\end{minipage} & \begin{minipage}[b]{\linewidth}\raggedright
Quote behavior
\end{minipage} \\
\midrule\noalign{}
\endhead
\bottomrule\noalign{}
\endlastfoot
Profile plus all required transparency evidence verified & Continue with
ownership-verified methods only \\
Transparency, checkpoint binding or operator continuity fails &
\texttt{no\_route} / explicit transparency error; never route \\
No resolver found a profile & \texttt{not\_registered} \\
Profile found but signature invalid & \texttt{invalid\_signature} \\
Profile expired & \texttt{no\_route} \\
Profile has no supported rail & \texttt{no\_route} \\
Transient resolver failures only & \texttt{not\_registered} unless a
stricter caller wants diagnostics \\
\end{longtable}

Implementations SHOULD expose resolver diagnostics in debug APIs, but
the user-facing quote response remains the stable protocol contract.

\hypertarget{conformance}{%
\subsection{Conformance}\label{conformance}}

A SatsPath resolver implementation is conformant if it:

\begin{itemize}
\tightlist
\item
  Returns signed profile data, not wallet secrets.
\item
  Does not decide payment execution.
\item
  Does not skip profile verification.
\item
  Distinguishes not found from invalid data.
\item
  Can be composed in a resolver chain.
\item
  Produces data compatible with the \texttt{QuoteResponse} pipeline.
\end{itemize}

\end{document}
